\section{Теоретический анализ правила выбора}\label{sec:theory}

\subsection{Ошибка сравнения кандидатов}
Зафиксируем признаки $X$, целевые метки $y$, пул разбиений $\mathcal P$ и разделение объектов на резерв $S$ и проверочную часть $T$. Обозначим $Q(P)=Q_T(P)$. Внешний лидер $\Pari$ максимизирует $A_L(P)$ на раскрытых метках $L\subseteq S$. Нас интересует, насколько ошибка такого сравнения влияет на качество выбранного разбиения.

Для кандидата $P$ определим расхождение доступной и проверочной оценок
\begin{equation}
 e_L(P)=A_L(P)-Q(P).
\end{equation}
При фиксированной схеме раскрытия меток это расхождение представляется как
\begin{equation}\label{eq:error-components}
 e_L(P)=b(P)+u_L(P),\qquad
 b(P)=\mathbb E_L[A_L(P)]-Q(P),\qquad
 \mathbb E_L[u_L(P)]=0.
\end{equation}
Величина $b(P)$ описывает систематическое расхождение, связанное со схемой получения меток, свойствами ARI и различиями между резервом и проверочной частью. Составляющая $u_L(P)$ характеризует вариативность по выборкам разметки. Равенство $b(P)=0$ для частичного ARI требует самостоятельного обоснования.

На выбор влияет порядок кандидатов. Поэтому введём максимальную ошибку парного сравнения
\begin{align}
 \rho_L
 &=\max_{P,P'\in\mathcal P}
 \left|\bigl(A_L(P)-A_L(P')\bigr)
       -\bigl(Q(P)-Q(P')\bigr)\right|\label{eq:pairwise-error}\\
 &=\max_{P\in\mathcal P}e_L(P)-\min_{P\in\mathcal P}e_L(P).\nonumber
\end{align}
Общий аддитивный сдвиг оценок сохраняет $\rho_L$. Эта величина учитывает именно искажение относительного качества кандидатов. Она доступна исследователю после раскрытия проверочных меток и служит диагностической характеристикой. Для $\epsilon_L=\max_P|e_L(P)|$ выполняется $\rho_L\le2\epsilon_L$.

\begin{proposition}[Потеря внешнего лидера]\label{prop:external}
Пусть $P^\star\in\arg\max_{P\in\mathcal P}Q(P)$. Тогда
\begin{equation}\label{eq:external-bound}
 R_T(\Pari)=Q(P^\star)-Q(\Pari)\le\rho_L\le2\epsilon_L.
\end{equation}
Если $|\mathcal P|\ge2$, оптимальное разбиение единственно и его отрыв
$\Delta=Q(P^\star)-\max_{P\ne P^\star}Q(P)$ удовлетворяет $\Delta>\rho_L$, внешний лидер совпадает с $P^\star$.
\end{proposition}
\begin{proof}
Максимальность $A_L(\Pari)$ даёт
\begin{align*}
 Q(P^\star)-Q(\Pari)
 &=A_L(P^\star)-A_L(\Pari)+e_L(\Pari)-e_L(P^\star)\\
 &\le e_L(\Pari)-e_L(P^\star)\le\rho_L.
\end{align*}
Для любого $P\ne P^\star$ из определения $\rho_L$ следует
$A_L(P^\star)-A_L(P)\ge\Delta-\rho_L>0$.
Следовательно, $P^\star$ также является единственным максимумом доступной оценки.
\end{proof}

Утверждение~\ref{prop:external} связывает качество выбора с точностью парных сравнений и разрывом между кандидатами. Оно справедливо для каждой реализованной выборки $L$. Вероятностная оценка правильности выбора потребовала бы дополнительно описать распределение $\rho_L$ при принятой схеме разметки. В эксперименте сопоставление $\rho_L$ и потери внешнего лидера позволяет оценить информативность доступных меток.

\subsection{Допуск и выбор по внутреннему индексу}
Введём $Q^\star=\max_{P\in\mathcal P}Q(P)$ и
$Q_c^\star=\max_{P\in\mathcal E_c}Q(P)$. Поскольку $\Pari\in\mathcal E_c$, допустимое множество содержит хотя бы одного кандидата. Добавляя и вычитая $Q_c^\star$, получаем точное разложение
\begin{equation}\label{eq:regret-decomposition}
 R_T(\Psilh)=
 \underbrace{Q^\star-Q_c^\star}_{D_{\mathrm{exclude}}}
 +\underbrace{Q_c^\star-Q(\Psilh)}_{D_{\mathrm{inside}}}.
\end{equation}
Первая составляющая измеряет потерю из-за исключения кандидатов на этапе допуска. Вторая измеряет потерю при выборе по внутреннему индексу внутри допустимого множества. Обе величины лежат в $[0,+\infty)$. Это разложение задаёт два направления диагностики: оценку состава допуска и оценку внутреннего ранжирования.

Потенциальное улучшение относительно внешнего лидера внутри допуска равно
\begin{equation}
 H_c=Q_c^\star-Q(\Pari)\ge0.
\end{equation}
Отсюда выигрыш окончательного выбора выражается как
\begin{equation}\label{eq:gain-decomposition}
 G_T=Q(\Psilh)-Q(\Pari)=H_c-D_{\mathrm{inside}}.
\end{equation}
Улучшение происходит при $H_c>D_{\mathrm{inside}}$. Таким образом, польза геометрии определяется одновременно наличием качественных альтернатив и способностью внутреннего индекса предпочесть их остальным допущенным разбиениям. Достаточное условие $D_{\mathrm{inside}}=0$ имеет вид
\begin{equation}
 \arg\max_{P\in\mathcal E_c}I(P)
 \subseteq\arg\max_{P\in\mathcal E_c}Q(P).
\end{equation}
Включение учитывает все варианты разрешения равенств внутренней оценки. При его выполнении $G_T=H_c\ge0$.

Для фиксированных бутстрепных перевыборок увеличение $c$ расширяет допустимое множество. Поэтому $D_{\mathrm{exclude}}$ убывает в широком смысле. Поведение $D_{\mathrm{inside}}$ зависит от вновь допущенных кандидатов и их внутренних оценок. Добавление кандидата с большим $I$ и меньшим $Q$ способно ухудшить окончательный выбор. Эта зависимость обосновывает отдельное исследование ширины допуска.

\begin{proposition}[Связь допуска с проверочным качеством]\label{prop:admission}
Для любого $P\in\mathcal E_c$ справедливо
\begin{equation}\label{eq:admission-bound}
 Q(\Pari)-Q(P)\le c\widehat\sigma_P+\varepsilon+\rho_L,
\end{equation}
где $\varepsilon$ --- численный допуск из~\eqref{eq:adaptive}. Для окончательного выбора также выполняется
\begin{equation}\label{eq:adaptive-bound}
 R_T(\Psilh)\le c\widehat\sigma_{\Psilh}+\varepsilon+\rho_L.
\end{equation}
\end{proposition}
\begin{proof}
Из определения $\rho_L$ и условия допуска получаем
\begin{align*}
 Q(\Pari)-Q(P)
 &\le A_L(\Pari)-A_L(P)+\rho_L\\
 &\le c\widehat\sigma_P+\varepsilon+\rho_L.
\end{align*}
Для доказательства~\eqref{eq:adaptive-bound} используем максимальность $A_L(\Pari)$:
\begin{align*}
 Q(P^\star)-Q(\Psilh)
 &\le A_L(P^\star)-A_L(\Psilh)+\rho_L\\
 &\le A_L(\Pari)-A_L(\Psilh)+\rho_L\\
 &\le c\widehat\sigma_{\Psilh}+\varepsilon+\rho_L.
\end{align*}
\end{proof}

Границы~\eqref{eq:admission-bound}--\eqref{eq:adaptive-bound} зависят от реализованной ошибки сравнения $\rho_L$. Их практическая информативность определяется величиной этой ошибки. Они показывают, что широкий допуск и существенное искажение внешнего ранжирования совместно допускают большие потери. Для оценки вероятности ухудшения требуется дополнительная связь между бутстрепной SD и ошибками парных сравнений.

\subsection{Частота переключений и их последствия}
При фиксированных данных и проверочной части случайность окончательного решения обусловлена выбором $L$ и конечным набором бутстрепных перевыборок. Вероятности и ожидания в этом подразделе учитывают оба источника. Индикатор $Z_L=\mathbf1\{\Psilh\ne\Pari\}$ отмечает смену разбиения. При $Z_L=0$ выигрыш равен нулю, поэтому для $\Pr(Z_L=1)>0$
\begin{equation}\label{eq:switch-effect}
 \mathbb E[G_T]=\Pr(Z_L=1)\mathbb E[G_T\mid Z_L=1].
\end{equation}
Условное среднее раскладывается на положительные и отрицательные эффекты:
\begin{equation}\label{eq:win-loss}
 \mathbb E[G_T\mid Z_L=1]=p_+\mu_+-p_-\mu_-.
\end{equation}
Здесь $p_+=\Pr(G_T>0\mid Z_L=1)$ и $p_-=\Pr(G_T<0\mid Z_L=1)$. При положительных вероятностях соответствующих событий
$\mu_+=\mathbb E[G_T\mid G_T>0,Z_L=1]$ и
$\mu_-=\mathbb E[-G_T\mid G_T<0,Z_L=1]$.
Для события нулевой вероятности соответствующий вклад в~\eqref{eq:win-loss} принимается равным нулю.

Разные разбиения могут иметь одинаковый проверочный ARI. Такие переключения дают нулевой вклад в средний эффект. Формулы~\eqref{eq:switch-effect}--\eqref{eq:win-loss} обосновывают совместное представление частоты переключений, вероятностей улучшения и ухудшения, а также их величин. Редкие крупные потери способны превысить суммарную пользу частых небольших улучшений.

\subsection{Роль парного бутстрепа}
Для двух заранее фиксированных кандидатов тождество дисперсий имеет вид
\begin{align}
 \operatorname{Var}_L\bigl(A_L(P_1)-A_L(P_2)\bigr)
 &=\operatorname{Var}_L(A_L(P_1))+\operatorname{Var}_L(A_L(P_2))\nonumber\\
 &\quad-2\operatorname{Cov}_L(A_L(P_1),A_L(P_2)).
\end{align}
Оценки кандидатов используют одни и те же размеченные объекты и могут сильно зависеть друг от друга. Парные перевыборки воспроизводят эту зависимость в условном бутстрепном распределении. Положительная ковариация уменьшает дисперсию разности относительно суммы маргинальных дисперсий.

В рассматриваемой процедуре $\Pari$ сначала выбирается на исходной разметке и затем фиксируется при бутстрепных сравнениях. Поэтому $\widehat\sigma_P$ описывает условный разброс разности при таком фиксированном лидере. Обоснование вероятностных свойств допуска требует учёта этапа выбора лидера~\citep{efron2014}, схемы раскрытия меток и вырожденных значений ARI. Систематическая составляющая $b(P)$ из~\eqref{eq:error-components} анализируется отдельно. При фиксированной разметке рост $B$ повышает точность вычисления бутстрепной SD относительно выбранного условного распределения.

\subsection{Покрытие классов и доступная информация}
Пусть два кандидата совпадают на размеченных объектах с точностью до перестановки номеров кластеров. Тогда их частичные ARI совпадают на исходной разметке и во всех объектных бутстрепных перевыборках. Отставание и SD парной разности равны нулю. Оба кандидата проходят адаптивный допуск, а выбор между ними определяется внутренним индексом и правилом разрешения равенств.

На скрытых объектах эти кандидаты могут по-разному объединять и разделять целевые классы. При наличии классов, полностью отсутствующих в $L$, доступная разметка характеризует лишь часть таких различий. Возможны разные продолжения целевых меток на скрытые объекты, согласующиеся с одной наблюдаемой разметкой и задающие противоположный порядок качества кандидатов. В этой ситуации предпочтение внутреннего индекса опирается на связь геометрии признаков с целевыми категориями. Именно эта связь определяет последствия выбора на проверочной части.

Полезность внутренней оценки поэтому следует исследовать совместно с покрытием классов. Случаи совпадения ограничений кандидатов на $L$ образуют отдельную диагностическую группу. Для них можно непосредственно проверить, как часто геометрическое разрешение равенств улучшает проверочный ARI и как величина эффекта зависит от состава наблюдаемых классов.

\subsection{Следствия для экспериментов}
Теоретический анализ задаёт три направления проверки. Первое --- оценка внешних сравнений через $\rho_L$, отрыв лучшего кандидата и фактическую потерю внешнего лидера. Второе --- раздельная диагностика $D_{\mathrm{exclude}}$ и $D_{\mathrm{inside}}$ при изменении порога допуска. Третье --- анализ частоты и последствий переключений по~\eqref{eq:switch-effect} и~\eqref{eq:win-loss}.

Повторные расчёты бутстрепа при фиксированном $L$ характеризуют вычислительную устойчивость допуска. Новые выборки $L$ характеризуют зависимость выбора от доступной разметки. Покрытие классов варьируется при одинаковом бюджете, а соответствие геометрии целевым меткам исследуется на независимых задачах. Такой протокол позволяет установить, в каких условиях внутренний индекс использует доступный потенциал улучшения и какие механизмы приводят к потерям.
